\chapter{Detection: the Positron}
\label{chap:detection-the-positron}

This is what the instrument has been built to reach. Through four parts the finite model has stood
beneath the physics on four kinds of honest floor --- a name, a count, a smooth-shadow analogy, a no-go wall
--- each modest in its own way. Here those floors give way to something firmer: a particle a finite apparatus
genuinely registers, by a logic the model proves. The detection is ours to build; the particle it catches is
the world's. The positron --- the electron's antiparticle, equal in mass and opposite in charge, the partner
the Dirac equation's negative-energy solutions predicted --- is the instrument's exemplar, the place where the
model proves the most, and, not by accident, the place where it fences the most clearly. The two go together:
the firmer the claim, the sharper the line around it must be drawn. Two effects carry the chapter: the
annihilation by which a positron is detected, and the threshold at which one is created.

\section{The Positron Annihilation Effect: a detection the model proves}
\label{se:positron-annihilation}

A positron is detected by its death. When a positron meets an electron the two opposite charges cancel and the
pair vanishes, its entire rest energy carried off as light --- two gamma photons,
\begin{equation}
  e^+ e^- \longrightarrow 2\gamma .
\end{equation}
Why two, and not one? A single photon cannot carry away the energy of a pair annihilating at rest: a lone
photon always has momentum, since it moves at the speed of light, but the pair at rest has none, and momentum
must balance. Two photons can balance it by departing in opposite directions, their momenta cancelling. So the
annihilation makes a pair of photons, and the energy they share is the two rest masses,
\begin{equation}
  E_\gamma = m_e c^2 = 511~\mathrm{keV},
\end{equation}
each photon bearing one electron's rest energy --- the 511~keV gamma line. Momentum conservation does more than
require two photons; it fixes their geometry. In the rest frame of the pair the total momentum is zero, so the
two photons must leave with equal and opposite momenta: equal energy, opposite directions, back-to-back along a
straight line through the point where they were born. The detector is ours to build; the two-photon signature
the world sends back is the world's.

In an ideal annihilation the pair is at rest, and the two photons are exactly 511~keV and exactly collinear. In
matter they are not quite. The electron a positron meets is bound in an atom and moving, so the pair carries a
small net momentum, and the books must still balance: that residual momentum tilts the two photons a fraction
of a degree off a straight line and shifts their shared energy a little above and below 511~keV. The line
broadens; the back-to-back angle opens slightly. The deviation is not noise --- it is a reading. The width of
the 511~keV line and the small departure from collinearity are the apparatus's measurement of the electron's
momentum in the material, the same conservation books read for the residue they leave. The annihilation
photon, looked at closely enough, carries off not only the rest energy but the motion of the electron that met
it --- so the line shape and the angular spread are a momentum spectrum of the bound electron, and have been
used to read the electron states of solids, mapping out where the electrons sit in a metal from the way their
motion blurs the annihilation line. The detector's residue is the material's portrait.

Often, too, the pair first binds for a fleeting instant before annihilating, into positronium --- a
hydrogen-like atom in which a positron stands in for the proton, electron and antiparticle circling their
common centre. It is a genuine atom of matter and antimatter, with its own ladder of levels close to
hydrogen's, and it comes in two spin arrangements: a singlet, whose two constituents annihilate into the two
back-to-back photons the detector is built to read, and a triplet, which annihilates instead into three
photons sharing the energy across a plane. Bound or free, singlet or triplet, the energy and momentum books
read the same; the coincidence apparatus is set for the two-photon line, and reads the positron by the death it
shares with an electron.

This geometry is what a coincidence detector reads. Set two detectors on opposite sides of a sample and watch
for two photons arriving together --- within a coincidence window of a few nanoseconds, both near the 511~keV
line, and back-to-back along the line joining the detectors. Three independent conditions: a time, an energy,
and a direction. When all three are met at once, a positron annihilated between them; when any one fails, the
apparatus is silent --- the three conditions ours to set, the coincidence the world fires the world's. The
registration is the conjunction, and the conjunction is checkable term by term.

Here the finite model does more than name or count or shadow --- it proves. The registration is a theorem. The
apparatus registers a positron if and only if both photons lie on the line and the pair is back-to-back,
\begin{equation}
  \text{registers} \;\Longleftrightarrow\; \text{on-line}(p) \,\wedge\, \text{on-line}(q) \,\wedge\,
  \text{back-to-back}(p, q),
\end{equation}
a genuine biconditional, not a one-way sufficiency. And the converses are real, falsifiable lemmas, each a way
the apparatus can be made to stay silent. A photon off the line is not registered: shift its energy out of the
window --- let it scatter once on the way out and arrive below 511~keV --- and the coincidence is rejected. A
pair that is not back-to-back is not registered: tilt the geometry, and the two detectors no longer face the
same line, so they no longer agree. A pair separated in time is not registered: let the two arrivals fall
outside the nanosecond window and the apparatus reads them as unrelated. Each converse is a separate dial the
experimenter can turn to force a silence, and each turning confirms one conjunct of the theorem. The
registration is logically equivalent to the conjunction of three checkable predicates, and the model proves the
equivalence --- a finite theorem about exactly when the apparatus fires. To this the model adds a proved
charge-cancellation it carries from its foundation: that the electron and positron records cancel, and that a
positron is what an asymmetric baseline yields. This is the instrument's strongest grounding.

And it is the most careful. The honest floor is a finite count obligation, the most fully proved the instrument reaches,
and its fence is drawn with equal precision --- because a claim this firm earns its force only if the line
around it is just as firm. What the model proves is the coincidence logic and the inherited charge-cancellation:
that, and exactly that. It does \emph{not} claim to detect real antimatter --- it proves when a coincidence
pattern registers and leaves the ontology of what causes it to physics. It does \emph{not} derive the 511~keV
line; that number is a labelled constant, the electron's rest energy fed in, not a result computed. It asserts
\emph{no} scattering amplitude --- nothing about how or why the annihilation proceeds, only the bookkeeping of
energy and momentum once it does. And it says \emph{nothing} about the rate: no cross-section, no probability
per encounter, no count per second. Each of these would be an overreach of a familiar kind: to derive the
511~keV would be to claim the model computes the electron's mass, when the mass is fed in; to assert an
amplitude or a rate would be to claim a dynamics the bookkeeping does not contain. The model proves the logic
of the coincidence and stops exactly there. The reading is that the positron is the exemplar not because the
model claims the most but because it proves the most while fencing the most clearly. The detection is a theorem,
converses and all; the charge-cancellation is proved, not posited; and in the same breath the model states
exactly what it leaves to physics --- the constant, the amplitude, the rate, the ontology. The strongest proof
the instrument reaches is also its clearest fence.

The positron was found exactly this way, by a track that bent the wrong way. Carl Anderson, in 1932,
photographed cosmic-ray tracks in a cloud chamber set in a magnetic field and caught one that curved as an
electron would --- but in the opposite sense, as though the electron carried a positive charge. The whole
inference hung on knowing which way the particle went, for a track curving one way for a downward positive
particle looks identical to one curving the other way for an upward negative particle. Anderson settled it with
a lead plate laid across the chamber: a particle crossing the plate loses energy and curves more tightly on the
far side, so the side of tighter curvature is the side it left by. The track he caught entered from below,
crossed the plate, and curved harder above --- moving upward, and bending as only a positive charge could. From
the curvature and the energy it shed in the lead he could read the mass as well: about an electron's, far too
light to be a proton. A positive particle of electron mass --- the positron, the first antiparticle seen,
confirming what Dirac's equation had predicted from its negative-energy solutions four years before.

And the annihilation is now put to work. Positron emission tomography --- a PET scan --- turns the exemplar's
logic into a clinical instrument. A tracer, most often a glucose analogue that collects where metabolism runs
highest, is injected; its nuclei emit positrons, each of which annihilates within a millimetre or so with a
nearby electron into two back-to-back 511~keV photons. A ring of detectors encircling the patient records each
coincident pair and draws the line between the two struck detectors --- the line of response --- along which
the annihilation must have lain --- the ring ours to build, the lines the annihilations write the world's.
Millions of such lines, intersected, reconstruct a three-dimensional map of
where the tracer collected; timing the two arrivals to fractions of a nanosecond narrows the annihilation to a
stretch of its own line and sharpens the picture further. The exemplar's coincidence logic, run in a hospital,
is a diagnosis.

The same line is written across the sky. From the direction of the galactic centre a steady glow at 511~keV has
been mapped by gamma-ray telescopes --- the light of positrons annihilating with electrons throughout the inner
galaxy, the exemplar's signal arriving from thousands of light-years away. Where all those positrons come from
is still argued --- radioactive nuclei forged in stars, or sources stranger --- but the line itself is
unmistakable, the electron's rest energy announced from across the galaxy. The coincidence the apparatus proves
on a bench and reads in a hospital, the sky has been emitting all along.

\section{The Positron Threshold Effect: a bracketed creation}
\label{se:positron-threshold}

Creating a positron, rather than detecting one, takes energy. Concretely it takes pair production: a photon
passing through the field of a nucleus converts into an electron--positron pair, possible only once the
photon's energy reaches twice the rest mass,
\begin{equation}
  E \ge 2 m_e c^2 = 1.022~\mathrm{MeV} .
\end{equation}
That a nucleus must be present is not incidental; it is conservation again, read the other way. A lone photon
converting to a pair in empty space would have to satisfy both books at once and cannot: in the frame where the
new pair is at rest the pair carries no momentum, but a photon always carries some, and nothing in the photon
and the pair alone can balance both energy and momentum. A third body must take up the slack. In pair
production that body is a nucleus, whose heavy field absorbs the recoil momentum at almost no cost in energy, so
that the photon's energy can pass almost entirely into the two rest masses. The threshold sits just above twice
the electron mass, the small excess paid to the recoiling nucleus. The same conservation that demanded two
photons in annihilation demands a third body in creation: the books balance, or the event does not occur ---
the energy ours to supply, the threshold the world sets on it.

Pair production is, in fact, the last of three ways a photon can spend itself in matter. At low energy a photon
is swallowed whole by an atom, ejecting a bound electron --- the photoelectric effect. At middling energy it
scatters off an electron and gives up part of its energy --- Compton scattering. Only above the 1.022~MeV
threshold does the third channel open, the photon converting into a pair; and at high energy it comes to
dominate the rest. The threshold is the energy at which a new way for light to disappear becomes available ---
the point above which a photon can turn into matter.

This is the Dirac negative-energy picture run forward. The Dirac equation's solutions came in two
signs, and to keep an electron from cascading endlessly down through the negative-energy states those states
are taken as already filled --- the sea --- with the exclusion principle barring any further electron from
joining them. Lift one electron out of the filled sea with enough energy and two things appear together: the
electron, now in a positive-energy state, and the hole it leaves behind, which carries itself in every way as a
particle of opposite charge --- the positron. Pair production is exactly that lift, and the threshold is the
energy it costs: the gap across the sea, twice the rest mass. The chirality and the Dirac operator described
the electron's two-signed dynamics; here the second sign is paid for and walks out
of the sea as a particle.

Below that energy a trial reads pure matter; at or above it a positron is admitted; and so there is a
threshold. A discrete apparatus cannot \emph{compute} that continuous threshold --- it has no access to the
smooth energy axis the threshold sits on --- but it can do something nearly as good: bracket it, and prove the
bracket holds a real crossing. Over the symmetric baseline the antimatter count is zero, and over the first
tilt it is at least one,
\begin{equation}
  \text{count}(\text{baseline}) = 0, \qquad \text{count}(\text{first tilt}) \ge 1,
\end{equation}
so the admitting threshold lies in the half-open interval $(0, 1]$, one unit wide on the apparatus's own scale
--- the tilt ours to make, the crossing it forces the world's. The crossing is not merely a change in a count; it is a change in the sign of the strain's second variation,
which the apparatus reads directly. At the flat baseline the second variation is negative, $\delta^2 = -1$, the
signature of a stable electron; at the first tilt it has flipped positive, $\delta^2 = +1$, the signature of a
positron-admitting configuration,
\begin{equation}
  \delta^2 : \; -1 \;\longrightarrow\; +1,
\end{equation}
a proved sign-change straddling the threshold. A sign-change is a sturdier thing to prove than a value: it
needs only the two ends, not the curve between, and the apparatus has the two ends. Where exactly inside the
bracket the crossing falls, the apparatus cannot say --- but absent any structure its records do not supply,
the unbiased estimate is the midpoint, the root of the straight line drawn between the two bracketing values,
\begin{equation}
  \text{linear-interpolation root} = \tfrac12,
\end{equation}
inside $(0, 1]$ and assuming nothing the records do not give.

The honest floor is a finite count obligation, and a proved one. The sign-change of the second variation and
the jump in the antimatter count from zero to at least one are proved; the exact threshold is bracketed to one
unit and estimated at $\tfrac12$, not computed. The reading is that a positron's creation is a proved threshold
the apparatus can bracket but not pinpoint. The count is zero below and one or more above, the second variation
flips from $-$ to $+$, and the model proves the crossing while honestly bracketing its location to one unit and
placing its unbiased estimate at the midpoint. A discrete apparatus cannot compute a continuous threshold, but
it can prove that one exists and bound where it lies.

\bigskip

With the positron, the instrument reaches its deepest grounding. Through naming, counting, analogy, and no-go the
finite model stood beneath the physics; here it stands level with it --- and is most careful exactly where it
proves the most. The annihilation's coincidence logic is a theorem with falsifiable converses, and the residual
momentum the photons carry off is read, not discarded; the threshold's sign-change and count-jump are proved,
the need for a third body is conservation read forward, and the location is honestly bracketed. The positron is
the exemplar because it is the one place where the finite apparatus does not approximate, shadow, or merely
name a phenomenon, but registers a particle by a proved logic --- the detecting ours, the particle it
registers the world's --- while stating, with the same precision, exactly what it does not claim. That double discipline, to prove all it can and fence all it cannot, is the
instrument's method in miniature, and the positron is where it shows clearest. Part IV closes next with decay,
confinement, and symmetry breaking.

\bigskip
\begin{center}
\rule{0.4\textwidth}{0.4pt}
\end{center}
\bigskip

\noindent Here the case reaches the exhibit it was built to secure, and it secures it the way it means to secure
everything: by proving all it can, and claiming not one inch beyond. The two go together, and the exemplar makes
it plain --- the firmer a claim, the sharper the line that must be drawn around it, or the force of it leaks out
into ground it never earned. So the instrument arrives at its strongest exhibit holding two disciplines at once:
to prove to the hilt, and to fence to the hilt, in the same breath.

The sign the last chapter proved leaves the page now and becomes a thing. What was a fact about loops --- a mark
a circuit carries home --- walks out of the theory as a particle the apparatus can genuinely catch: the
electron's antiparticle, its equal and its opposite, alike in mass and reversed in charge. And the apparatus
catches it by its death. Let it meet an electron and the two opposite charges cancel, the pair vanishing into
light --- a matched pair of photons flung back to back, equal and opposite, along a line through the point where
the two died. That paired flash is the signature the detector is built to read: the particle registered not by
being named, or counted, or shadowed, but by a real event the record can point to. It was first caught much this
way, by a track that bent the wrong way --- a particle in a cloud chamber curving as an electron would but in
the opposite sense, and a plate of lead laid across to settle which direction it truly ran. A positive particle
of an electron's slight mass: the first antiparticle ever seen, exactly the partner the theory had predicted
from its equations four years before it was found.

And here the instrument does the strongest thing in the book, and --- not by accident --- the most careful. The
detector fires on a coincidence of three plain conditions: two photons arriving together, both at the right
energy, and back to back along the line joining the detectors. It registers the particle if and only if all
three hold at once, and the \emph{if and only if} is the whole of the strength. It is a theorem, converses and
all: turn any one dial and force a silence --- soften a photon's energy below the line, tilt the pair off the
straight, let the two arrivals drift apart in time --- and each silence confirms one clause of the proof. The
instrument does not merely say the coincidence suffices; it says it is exactly what registration \emph{is}, both
directions in hand. This is the place where the finite model stops standing beneath the physics and stands level
with it. And in the very same breath it draws the line around what it has shown. It proves the coincidence logic
and the cancellation of the two charges --- that, and precisely that. It does not claim to have caught
antimatter itself, only to have proved when the pattern fires, and leaves what causes it to physics; it does not
derive the energy of the flash, which is a constant fed in and not a number it computed; it asserts nothing
about how the annihilation proceeds, nor how often it does. Each of those would be an inch more than the theorem
gives, and the instrument takes not one of them. Its strongest proof is, in the same stroke, its clearest fence.

The same double discipline governs the making of one, not only the catching. To create the particle takes
energy, and a third body to carry off the recoil, or the books of energy and momentum cannot both close and the
event does not happen at all. The exact energy at which the crossing comes the apparatus cannot compute --- the
smooth axis it would need is not one it holds --- but it does the honest thing in its place: it brackets the
crossing between a baseline that admits none and a first step that admits at least one, proves that a real
crossing lies between them, and, claiming nothing its records do not give, sets its estimate at the midpoint. It
proves that a threshold is there, and fences, to the width of a single step, exactly how far it can say where.

And the \emph{charge} comes home here in the plainest sense of all. The charge that was a count, and then a
signed count, is now a particle --- a thing of opposite charge, caught and registered and countable on a
detector, the sign of the last chapter made flesh and pointed at. More than that: the way it is caught is a
cancellation of charge against charge, the positive against the negative, the two annihilating into light --- so
the detector reads the charge in the very act of its cancelling. What the case will collect on is no longer an
abstraction; it is a thing that leaves a track, bends in a field, and can be shown. And the logic that catches
it has long since left the bench: run in a hospital, the same coincidence draws its lines through a patient and
maps where a tracer gathered; turned on the sky, it reads a steady glow of these very deaths from the direction
of the galaxy's heart, thousands of light-years off. The exemplar's proof is a diagnosis and a beacon at once. And the trace is the paired
flash that outlives the death --- the two photons the record keeps as the mark of a charge that was there and is
now spent. Whatever comes to be owed rides on a charge the detector can point to.

A registered particle, proved by a theorem, is the surest single reading the case will ever enter. But it is
still a single reading, and the case is not made of one. The wheel's account, the field's, and the one still to
come from overhead must be brought to agree as they always must; a proof, even this one, binds the thing it
proves and not the three hands that must still meet. The exemplar is the firmest exhibit in the file --- and the
file is still open.

With this the instrument has reached the exhibit it was built for, and Part IV has one turn left. The next
chapter follows the particle into its undoing --- how it decays, how it is held confined, how a symmetry the
world once kept comes to be broken --- and closes the quantum part. The case is still open, its firmest exhibit
entered and still unmet by the others; the court does not sit until the record is complete.
